% =============================================================================
% CoLM 2026 — Submission 991: "How Reliable Are VLM Judges? Conformal Prediction
%                              Reveals Task-Dependent Uncertainty in Multimodal
%                              Evaluation"
%
% This file contains the THREE official reviewer reviews verbatim (typos and
% formatting preserved as posted on OpenReview). No rebuttal content.
% =============================================================================
\documentclass[10pt]{article}

\usepackage[margin=0.7in,top=0.6in,bottom=0.7in]{geometry}
\usepackage{amsmath, amssymb}
\usepackage{xcolor}
\usepackage{hyperref}
\usepackage{enumitem}
\usepackage{parskip}
\usepackage[compact]{titlesec}
\usepackage{mathptmx} % Times Roman — much denser & very readable
\usepackage[T1]{fontenc}
\hypersetup{colorlinks=true, linkcolor=blue, urlcolor=blue, citecolor=blue}

% Tighten section spacing
\titlespacing*{\section}{0pt}{0.6ex plus .2ex minus .2ex}{0.3ex plus .1ex}
\titlespacing*{\subsection}{0pt}{0.4ex plus .15ex minus .15ex}{0.15ex plus .05ex}
\titleformat{\section}{\normalfont\large\bfseries}{\thesection}{0.6em}{}
\titleformat{\subsection}{\normalfont\normalsize\bfseries}{\thesubsection}{0.5em}{}
\titleformat{\paragraph}[runin]{\normalfont\normalsize\bfseries}{}{0pt}{}[\,]

% Denser reviewer quote: smaller indent, no extra vertical space, footnote size
\newenvironment{reviewerquote}{%
  \par\addvspace{0.25ex}\begingroup
  \footnotesize\itshape
  \leftskip=0.6em \rightskip=0.6em
  \color{black!78}
  \setlength{\parskip}{0.35ex plus 0.1ex}
}{%
  \endgroup\par\addvspace{0.25ex}
}

% Compact lists
\setlist{nosep,leftmargin=*,topsep=0.15ex,partopsep=0pt}

\title{\vspace{-1.5em}CoLM 2026 --- Submission 991\\[0.2em]
       \normalsize Reviews from Three Reviewers (Verbatim from OpenReview)}
\author{}
\date{}

\begin{document}
\maketitle
\vspace{-2.5em}

\noindent\textbf{Score summary.}\quad
GBAF: \textbf{7} (Good paper, accept) \quad$\bullet$\quad
WoGp: \textbf{5} (Marginally below threshold) \quad$\bullet$\quad
YU66: \textbf{4} (Rejection) \quad$\bullet$\quad
Average: \textbf{5.33}

% =============================================================================
\section{Reviewer GBAF (Rating: 7 --- Accept)}
% =============================================================================

\paragraph{Title of review.} \emph{Conformal Prediction to Better Evaluate of VLM Judge Reliability}

\subsection*{Summary}
\begin{reviewerquote}
This work explores the reliability of VLM judges, particularly diving into correlation with humans on ranking tasks does not equal reliability on point judgements. To formalize this, the authors use conformal prediction, where point scores are converted into conformal intervals. This gives the uncertainty around any given model judgement. They use two multimodal datasets with human-annotated scores for calibration.

They find that there is a disparity between \emph{ranking}, which is consistent amongst judges, and \emph{scoring}, which is unreliable with variance. This distinction is seen in particular across different tasks and is dependent on the data annotation quality. They suggest using conformal preidction as a means of determining the best evaluation setting: when intervals are narrow, absolute scoring is appropriate; when intervals are wide, relative ranking is more reliable.
\end{reviewerquote}

\subsection*{Reasons To Accept}
\begin{reviewerquote}
\textbf{well grounded motivation:} I find this to be a really interesting motivation with informative results! VLM judges are being used widely to automate evaluation and correlation with humans is the most common way to validate the approach. This works tackles the importance of exploring \emph{reliability}, beyond just high correlation numbers.

\textbf{strong findings about correlation/accuracy versus reliability:} As noted in the paper, we tend to rely on correlation with human ratings to justify automatic evaluations using judges. The findings in this work, particularly in Section 5.2, really add depth about how high performing models wrt ranking can still have very poor point score estimations.

\textbf{depth of evaluation:} The authors tested a range of different conformal prediction methods (Table 1). See questions section for additional take on this contribution.

\textbf{range of VLM judges used:} While only 3 VLMs are evaluated, I appreciate that these were intentionally chosen with one ``speciality'' judge model, one slightly larger reasoning model and one higher performing (on general multimodal tasks) closed-source model.
\end{reviewerquote}

\subsection*{Reasons To Reject}
\begin{reviewerquote}
\textbf{limited comparison of tasks:} I appreciate the number of tasks evaluated in Table 3. However, the categories feel a bit arbitrary --- COCO is ``Gen. VQA'' yet TextVQA is ``Vision''. To draw more structured conclusions about the task-dependent differences, I suggest making a more formal taxonomy to define the tasks, perhaps drawing on existing work.
\end{reviewerquote}

\subsection*{Questions To Authors}
\begin{reviewerquote}
\textbf{limited details on different conformal methods:} Unless readers are already quite familiar with conformal prediction methods, the names in Section 5.1 might be unfamiliar. I suggest adding a citation to each. As a minor note, this contributes less to the overall narrative and could probably be moved to the conclusion.
\end{reviewerquote}

\paragraph{Rating.} 7: Good paper, accept. \quad \textbf{Confidence:} 3 (fairly confident). \quad \textbf{Ethics Flag:} No.

% =============================================================================
\section{Reviewer WoGp (Rating: 5 --- Marginally Below Acceptance Threshold)}
% =============================================================================

\paragraph{Title of review.} \emph{a useful empirical study of conformal intervals for VLM judges, but unclear conformal validity}

\subsection*{Summary}
\begin{reviewerquote}
The paper studies whether VLM-as-a-Judge scores can be accompanied by calibrated prediction intervals using conformal prediction over score-token log-probabilities, evaluated on MLLM-as-a-Judge and Polaris with LLaVA-Critic-7B, Phi-4-reasoning-vision-15B, and Gemini 2.5 Flash. The work is timely and potentially significant: Tables 1--3 show that point-score metrics hide substantial task-dependent uncertainty, and the ranking-scoring decoupling analysis in \S 5.5 is a useful warning against relying only on Pearson/Spearman correlations. However, the paper's methodological contribution is incremental relative to existing conformal LLM-as-a-judge work, which already introduced interval evaluation, ordinal boundary adjustment, and midpoint analysis for text-only judges; the novelty is mainly the multimodal empirical stratification rather than a new conformal framework. More importantly, the paper has unresolved issues around the validity of the conformal protocol, inconsistent descriptions of boundary adjustment, and several numerical inconsistencies in interval-width normalization, which make the main reliability claims harder to trust in their current form.
\end{reviewerquote}

\subsection*{Reasons To Accept}
\begin{reviewerquote}
The paper targets an important failure mode in multimodal evaluation: Table 1 shows that all three VLM judges have low exact agreement with human scores, around 32--34\%, and modest correlations, with Pearson ($\rho$) ranging from .303 to .459, so adding uncertainty estimates is practically motivated rather than cosmetic.

The empirical scope is stronger than a single-model case study: \S 4 evaluates three heterogeneous judges, two datasets, eight conformal methods, and 10 random splits, while Tables 2 and 9 provide a useful comparison showing that R2CCP and CHR dominate degenerate alternatives such as CQR and OrdinalAPS, which often collapse to nearly full-range intervals.

The task-level analysis is the paper's most compelling contribution: Table 3 and Figure 1 show that raw R2CCP widths vary from 2.08 on AesBench to 3.50 on InfographicsVQA for LLaVA-Critic, and Table 6 shows wider intervals for Vision-Heavy tasks than Knowledge/Web tasks, supporting the central claim that VLM-judge reliability is task-dependent.
\end{reviewerquote}

\subsection*{Reasons To Reject}
\begin{reviewerquote}
The conformal prediction protocol is not specified clearly enough to justify the claimed coverage guarantees: \S 3.2 says the point predictor $f(x)$ is ``trained on the calibration set,'' and \S 4 says the data are split only 50/50 into calibration and test sets, but standard split conformal validity requires that the model used to produce nonconformity scores be fixed before calibration, or that training/calibration be separated; if R2CCP or other predictors are trained and calibrated on the same examples, the guarantee in Equation 4 does not directly apply.

The boundary-adjustment description appears internally inconsistent with the reported results: \S 3.3 states that the interval $[l,u]$ is mapped to $[\lceil l\rceil, \lfloor u\rfloor]$, which should weakly shrink intervals on an integer rating scale, yet Tables 1 and 2 report boundary-adjusted intervals that are much wider than raw intervals, e.g.\ LLaVA-Critic R2CCP width increases from 3.05 to 3.60 and coverage increases from .900 to .981; this suggests either the formula, implementation, or explanation is wrong.

Several quantitative claims about interval width normalization are inconsistent: \S 4 defines width on a 1--5 scale as ranging from 0 to 4, and \S 5.4 correctly describes AesBench width 2.08 as 52\% of the score range and InfographicsVQA width 3.50 as 88\%, but Table 7 describes width 3.05 as 61\% rather than $3.05/4 \approx 76\%$, and the abstract's ``$\sim$40\%'' figure for aesthetics/natural images is not supported by Table 3.

The claim that interval width is ``primarily driven by task difficulty and annotation quality'' is over-causal relative to the evidence: the Polaris comparison in \S 5.6 and Table 7 changes multiple factors simultaneously, including task type, score aggregation, number of annotators, label continuity, score distribution, and benchmark construction, so the 4.5$\times$ width reduction cannot isolate annotation quality without controlled ablations.

The ranking-scoring gap is interesting but underformalized: Equation 7 is missing necessary parentheses, since the values in Table 17 appear to use $\mathrm{RSG}(d) = |\rho_d| - (1 - w_d/(K-1))$, not the printed $(|\rho_d| - 1 - w_d/(K-1))$; moreover, the metric is not validated against downstream decisions and can become positive simply because intervals are very wide, so it should be presented as an exploratory diagnostic rather than a ``fundamental'' failure mode.
\end{reviewerquote}

\paragraph{Rating.} 5: Marginally below acceptance threshold. \quad \textbf{Confidence:} 3 (fairly confident). \quad \textbf{Ethics Flag:} No.

% =============================================================================
\section{Reviewer YU66 (Rating: 4 --- Rejection)}
% =============================================================================

\paragraph{Title of review.} \emph{An empirical study of the application of conformal prediction to VLM-as-a-Judge found interesting but limited insights}

\subsection*{Summary}
\begin{reviewerquote}
This work adapts conformal prediction to the VLM-as-a-Judge setting. It constructs a five-dimensional feature vector from score-token log-probabilities and evaluates eight conformal methods on two benchmarks: MLLM-as-a-Judge (14 single-annotation tasks) and Polaris (captioning with multiple annotators). The study compares three judge models, LLaVA-Critic-7B, Phi-4-reasoning-vision-15B, and Gemini 2.5 Flash.

Key findings include: R2CCP achieves the strongest overall performance on a five-point Likert scale with around 90\% coverage and an average raw interval width of approximately 3.05. Interval width varies substantially across tasks, with narrower intervals on AesBench (2.08) and wider ones on InfographicsVQA (3.50), suggesting that visually intensive reasoning tasks induce greater uncertainty. The study reports a ranking--scoring decoupling phenomenon, where high correlation does not imply tight or reliable score intervals, as illustrated on ChartQA. Intervals shrink by a factor of 4.5 on Polaris relative to MLLM-as-a-Judge. Mondrian CP reduces interval width by about 16.6\% on easier tasks. Feature fusion across multiple judges degrades performance.
\end{reviewerquote}

\subsection*{Reasons To Accept}
\begin{reviewerquote}
The use of VLMs as evaluators has become widespread, yet systematic reliability quantification remains limited. Introducing conformal prediction into this setting provides a useful perspective for the community.

The empirical study spans multiple judge models, a diverse set of 14 task categories, and several conformal methods, offering a reasonably comprehensive experimental picture. Observations such as task-dependent interval width, the tendency of judges to over-score low-quality answers and under-score high-quality ones, and the variability across benchmarks are practically informative.

The ranking--scoring decoupling phenomenon is worth highlighting. Strong ordinal correlation does not guarantee reliable absolute scores, which has direct implications for model selection, data filtering, and automated benchmarking pipelines based on VLM judges.
\end{reviewerquote}

\subsection*{Reasons To Reject}
\begin{reviewerquote}
The conformal protocol is insufficiently specified and may invalidate the claimed coverage guarantees. The paper states that $f(x)$ is trained on the calibration set while the same set is also used to compute nonconformity quantiles. For methods such as R2CCP, CQR, and CHR that require learning conditional distributions or prediction intervals, the absence of a proper training/calibration split, cross-conformal procedure, or leave-one-out scheme breaks the theoretical guarantees. The paper should clearly describe the data partitioning strategy and report results under a valid protocol.

The prediction intervals are excessively wide, which limits their practical usefulness. R2CCP yields a raw width of about 3.05 out of 4 and a boundary-adjusted width of about 3.60 out of 4, with several methods approaching the full range ([1,5]) after adjustment. While coverage is achieved, the intervals provide little actionable information at the instance level. This limitation is not sufficiently examined.

The narrower intervals observed on Polaris admit multiple explanations, including averaging over annotators, the homogeneous captioning task, smoother score distributions, better alignment between judges and task, lower label noise, or more concentrated targets. A comparison between MLLM-as-a-Judge and Polaris alone does not support the claim that interval width is primarily driven by annotation quality and task structure.

The ranking-scoring gap is introduced but not systematically analyzed. The paper does not report the metric with statistical rigor such as confidence intervals or significance tests, and relies mainly on illustrative examples like ChartQA and WIT. As a proposed failure mode, the treatment remains descriptive. The decoupling between ordinal correlation and absolute accuracy has been discussed in prior work on rating models, recommender systems, and LLM-based evaluation, including issues such as verbosity and position bias or calibration breakdown under distribution shift. The claim that this is a previously unrecognized failure mode is overstated. The definition in Equation (7) is also ad hoc, combining $|\rho|$ and $1 - w/(K-1)$ without a clear theoretical basis or connection to established calibration--discrimination decompositions such as Brier decomposition or expected calibration error.
\end{reviewerquote}

\paragraph{Rating.} 4: Ok but not good enough --- rejection. \quad \textbf{Confidence:} 3 (fairly confident). \quad \textbf{Ethics Flag:} No.

\end{document}
